% =================================================================
% BÖLÜM 2 — SİSTEM MİMARİSİ
% =================================================================
\section{Sistem Mimarisi}

% ---------------------------------------------------------------
\begin{frame}{5 Katmanlı Yazılım Mimarisi}{Tek yönlü bağımlılık ilkesi}
  \begin{columns}[c]
    \begin{column}{0.62\textwidth}
      \centering
      \includegraphics[width=\linewidth,height=0.72\textheight,keepaspectratio]{layered_arch.png}
    \end{column}
    \begin{column}{0.38\textwidth}
      \scriptsize
      Her katman yalnızca altındakiyle konuşur; donanım bağımlılığı izole edilir.
      \begin{itemize}\setlength{\itemsep}{2pt}
        \item \textbf{Uygulama} — imu / baro / gnss / telemetry task
        \item \textbf{Algoritma} — Mahony, Kalman, FSM, SUT
        \item \textbf{Sürücü} — bmi088, bme280, NMEA parser
        \item \textbf{HAL + Middleware} — CubeMX, FreeRTOS, SEGGER
        \item \textbf{Donanım} — BMI088, BME280, L86, USART
      \end{itemize}
      \vspace{2pt}
      \begin{block}{\footnotesize Modülerlik}
        \scriptsize SIT/SUT test katmanları aynı mimariye temiz biçimde eklenir.
      \end{block}
    \end{column}
  \end{columns}
\end{frame}

% ---------------------------------------------------------------
\begin{frame}{Donanım Haritası}{STM32F446 — Cortex-M4 @ 180 MHz, FPU, 512 KB Flash / 128 KB SRAM}
  \centering
  \includegraphics[width=0.92\textwidth,height=0.48\textheight,keepaspectratio]{hw_block.png}
  \par\vspace{3pt}
  \begin{columns}[T]
    \begin{column}{0.42\textwidth}
      \centering\scriptsize
      \begin{tabular}{@{}ll@{}}
        \toprule
        \textbf{Birim} & \textbf{Görev} \\
        \midrule
        I2C1    & BMI088 IMU \\
        I2C3    & BME280 Baro \\
        EXTI3/4 & BMI088 DRDY Kesmesi \\
        USART6  & GNSS (L86) \\
        USART2  & Telemetri / SIT \\
        IWDG    & Donanım Watchdog \\
        \bottomrule
      \end{tabular}
    \end{column}
    \begin{column}{0.58\textwidth}
      \begin{block}{}
        \scriptsize
        \begin{itemize}
          \item \textbf{DMA} — CPU'yu meşgul etmeden veri taşır
          \item \textbf{NVIC + EXTI} — asenkron veri-hazır kesmeleri
          \item \textbf{FPU} — quaternion float işlemleri donanımsal olarak hızlı\\
            \tiny \textit{FPU (Floating Point Unit): kayan nokta işlemlerini yazılımsal\\
            emülasyon yerine donanımda yapar; Mahony/Kalman hesapları\\
            CPU'yu daha az meşgul eder.}
        \end{itemize}
      \end{block}
    \end{column}
  \end{columns}
\end{frame}

% ---------------------------------------------------------------
\begin{frame}{FreeRTOS Görev Modeli}{Sabit öncelikli, birbirini kesebilen dört görev}
  \begin{table}
    \centering\small
    \begin{tabular}{@{}llll@{}}
      \toprule
      \textbf{Görev} & \textbf{Öncelik} & \textbf{Tetikleyici} & \textbf{İşlem / Periyot} \\
      \midrule
      \texttt{imu\_task}       & High        & Donanım kesmesi    & Mahony + snapshot \\
      \texttt{baro\_task}      & BelowNormal & 100 ms periyot     & Kalman + Flight SM \\
      \texttt{gnss\_task}      & BelowNormal & DMA Circular RX    & 1 Hz NMEA parse \\
      \texttt{telemetry\_task} & BelowNormal & İletişim hazır     & 50 Hz UART TX \\
      \bottomrule
    \end{tabular}
  \end{table}
  \vspace{6pt}
  \begin{exampleblock}{Deterministik veri zinciri (IMU)}
    \centering\small
    DRDY Pin $\rightarrow$ EXTI ISR $\rightarrow$ \texttt{xTaskNotifyFromISR} $\rightarrow$
    Görev uyanır $\rightarrow$ DMA başlat $\rightarrow$ DMA-Complete IRQ $\rightarrow$
    Parse $\rightarrow$ Snapshot
  \end{exampleblock}
  \centering\footnotesize
  ISR içinde iş minimumda: sadece bildirim gönderilir, hesaplama göreve devredilir.
\end{frame}

% ---------------------------------------------------------------
\begin{frame}[fragile]{Senkronizasyon Kalıpları}{Bekleme süresini minimuma indiren iki desen}
  \begin{columns}[T]
    \begin{column}{0.5\textwidth}
      \centering
      \textbf{\small ISR $\rightarrow$ Görev: Task Notification}\\[2pt]
      \includegraphics[width=\linewidth,height=0.38\textheight,keepaspectratio]{task_notify.png}
\begin{lstlisting}
void HAL_I2C_MemRxCpltCallback(...)
{
  BaseType_t woken = pdFALSE;
  xTaskNotifyFromISR(s_handle,
      NOTIFY_DMA_DONE, eSetBits, &woken);
  portYIELD_FROM_ISR(woken);
}
\end{lstlisting}
    \end{column}
    \begin{column}{0.5\textwidth}
      \centering
      \textbf{\small Görev $\rightarrow$ Görev: Snapshot Queue}\\[2pt]
      \includegraphics[width=\linewidth,height=0.44\textheight,keepaspectratio]{snapshot_queue.png}
      \vspace{2pt}
      \begin{block}{\footnotesize Derinliği 1 kuyruk (mailbox)}
        \scriptsize
        \begin{itemize}
          \item \texttt{xQueueOverwrite} — üretici hep son veriyi yazar
          \item \texttt{xQueuePeek} — tüketici eksiltmeden okur
          \item Tek sahipli, atomik, gereksiz tamponlama yok
        \end{itemize}
      \end{block}
    \end{column}
  \end{columns}
\end{frame}
